\documentclass[blue]{beamer}
\usetheme{metropolis}
\usepackage{amssymb,latexsym}
\usepackage{amsmath}
\usepackage{amsthm}
\usepackage{graphicx}
\usepackage{subfigure}
\usepackage{hyperref}
\usepackage{subcaption}
\definecolor{Red}{rgb}{1,0,0}
\definecolor{Blue}{rgb}{0,0,1}
\definecolor{darkblue}{rgb}{0,0,.75}
\definecolor{Green}{rgb}{0,1,0}
\definecolor{magenta}{rgb}{1,0,.6}
\definecolor{lightblue}{rgb}{0,.5,1}
\definecolor{lightpurple}{rgb}{.6,.4,1}
\definecolor{gold}{rgb}{.6,.5,0}
\definecolor{orange}{rgb}{1,0.4,0}
\definecolor{hotpink}{rgb}{1,0,0.5}
\definecolor{newcolor2}{rgb}{.5,.3,.5}
\definecolor{newcolor}{rgb}{0,.3,1}
\definecolor{newcolor3}{rgb}{1,0,.35}
\definecolor{darkgreen1}{rgb}{0, .35, 0}
\definecolor{darkgreen}{rgb}{0, .6, 0}
\definecolor{darkred}{rgb}{.75,0,0}
%\usepackage{beamerthemesplit}
\usepackage{color}
\usepackage{multirow}

\usepackage{lipsum}
\usefonttheme{professionalfonts}
\newcommand\blfootnote[1]{%
  \begingroup
  \renewcommand\thefootnote{}\footnote{#1}%
  \addtocounter{footnote}{-1}%
  \endgroup
}
\newcommand{\E}{\mathbb{E}}
\newcommand{\bi}{\begin{itemize}}
\newcommand{\ei}{\end{itemize}}
%\AtBeginSection{}
\setbeamertemplate{section in toc}[sections numbered]
\setbeamerfont{itemize/enumerate body}{size=\normalsize}
\setbeamerfont{itemize/enumerate subbody}{size=\normalsize}

%\usepackage{amsmath,amssymb,amsthm}
%\usepackage{graphicx}
\usepackage{booktabs}
\usepackage{tikz}
\usetikzlibrary{arrows.meta,positioning,decorations.pathreplacing}
%\usepackage{hyperref}
\title[Chapter3]{Chapter 3: The Solow Model}
\author[]{}

\date[May 2026]{May 2026}





\begin{document}
  \frame
  {
    \titlepage
  }


% ============================================================
% OUTLINE
% ============================================================
\begin{frame}{Outline}
	\tableofcontents
\end{frame}
% ============================================================
\section{Introduction}
% ============================================================

\begin{frame}{Motivation}
	\textbf{Motivating fact}: The capital-output ratio $K_t/Y_t$ is roughly stable at $\approx 3$ over time.
	\begin{itemize}
		\item Earlier theories (Harrod-Domar): treated this as a technological property (rigid, fixed proportions)
		\item But capital and labor are substitutable
		\item Why is $K_t/Y_t$ approximately constant despite substitutability?
	\end{itemize}
	
	\textbf{Solow's insight}: Reconcile the stable capital-output ratio  through \textbf{neoclassical} forces (decreasing marginal returns).
	

	\textbf{Two main takeaways}:
	\vspace{-6pt}
	\begin{enumerate}
		\item The fundamental source of long-run per capita income growth is growth in technology $A_t$
		\item If parameter values are common, different economies \textbf{converge} to the same income per capita in the long run
	\end{enumerate}
\end{frame}

% ============================================================
\section{The Basic Model}
% ============================================================

\begin{frame}{The Aggregate Production Function}
	\[
	Y_t = F(K_t, L_t)
	\]
	
	\textbf{Assumptions} on $F(K,L)$:
	\begin{enumerate}
		\item Strictly increasing in both $K$ and $L$
		\item Strictly quasiconcave (strictly convex isoquants)
		\item Constant returns to scale: $F(cK, cL) = cF(K,L)$ for all $c > 0$
		\item $F(0, L) = 0$
		\item Inada conditions: $\lim_{K \to 0} F_1(K,L) = \infty$ and $\lim_{K \to \infty} F_1(K,L) = 0$
	\end{enumerate}
	
	\bigskip
	Normalize population and hours to 1: $L_t = 1$. Define $f(k_t) \equiv F(k_t, 1)$:
	\[
	y_t = f(k_t)
	\]
\end{frame}

\begin{frame}{Capital Accumulation and the Fundamental Equation}
	\textbf{Capital accumulation}:
	\[
	k_{t+1} = i_t + (1-\delta)k_t
	\]
	
	\textbf{Goods market clearing}: $y_t = c_t + i_t$
	
	\textbf{Saving behavior} (exogenous, constant saving rate):
	\[
	i_t = sy_t, \qquad s \in (0,1)
	\]
	
	\bigskip
	Substituting:
	\[
	k_{t+1} = (1-\delta)k_t + sf(k_t)
	\]
	
	This is the \textbf{fundamental equation of the Solow model}. Given $k_0$, the entire path $\{k_t, y_t, c_t, i_t\}$ is determined.
\end{frame}

\begin{frame}{Steady State}
	A \textbf{steady state} is a constant $\bar{k}$ satisfying $k_t = \bar{k}$ for all $t$:
	\[
	\bar{k} = (1-\delta)\bar{k} + sf(\bar{k}) \qquad \Rightarrow \qquad \delta \bar{k} = sf(\bar{k})
	\]
	
	\begin{itemize}
		\item LHS: straight line through origin with slope $\delta$
		\item RHS: strictly increasing, strictly concave, slope $\infty$ at origin, slope $\to 0$
		\item $\Rightarrow$ Unique positive intersection (guaranteed by Inada conditions)
	\end{itemize}
	
	\bigskip
	\textbf{Cobb-Douglas}: $f(k) = k^\alpha$, $\alpha \in (0,1)$. Closed-form:
	\[
	\bar{k} = \left(\frac{s}{\delta}\right)^{\frac{1}{1-\alpha}}
	\]
	
	Steady-state $k/y$ ratio: $\bar{k}/\bar{y} = s/\delta$.
\end{frame}

\begin{frame}{Dynamics and Convergence}
	Dynamics in the Solow model
				\vspace{-5pt}
	\begin{figure}[h]
		\centering
		\includegraphics[scale=0.42]{FigsChapter3/Solow1R}
	\end{figure}
	\vspace{-13pt}
	Starting from $k_0$: use the curve to get $k_1$, place $k_1$ on horizontal axis, repeat.
	
	\textbf{Global, monotonic convergence} to $\bar{k}$ regardless of $k_0$:
		\vspace{-6pt}
	\begin{itemize}
		\item If $k_t < \bar{k}$: $sf(k_t) > \delta k_t$ $\Rightarrow$ $k_{t+1} > k_t$ (capital rising)
		\item If $k_t > \bar{k}$: $sf(k_t) < \delta k_t$ $\Rightarrow$ $k_{t+1} < k_t$ (capital falling)
	\end{itemize}
\end{frame}


\begin{frame}{Intuition for Convergence}
	Rewrite the capital accumulation in growth-rate form:
	\[
	\frac{k_{t+1} - k_t}{k_t} = s\frac{f(k_t)}{k_t} - \delta
	\]
	
	\begin{itemize}
		\item $f(k)/k$ is \textbf{decreasing} in $k$ (from $f''(k) < 0$ and $f(0) = 0$)
		\item When $k$ is small: output per unit of capital is large $\Rightarrow$ high investment relative to capital $\Rightarrow$ capital grows fast
		\item When $k$ is large: output per unit of capital is small $\Rightarrow$ depreciation dominates $\Rightarrow$ capital falls
		\item \textbf{Key}: decreasing returns to capital generate the convergence force
	\end{itemize}
	
	\bigskip
	This de-mystifies the stability of $k/y$: deviations from $\bar{k}/\bar{y} = s/\delta$ are self-correcting.
\end{frame}

\begin{frame}{Other Kinds of Dynamics: Endogenous Growth}
	Endogenous growth in the Solow model
			\vspace{-5pt}
\begin{figure}[h]
	\centering
	\includegraphics[scale=0.42]{FigsChapter3/Solow3R}
\end{figure}
\vspace{-13pt}	
	\textbf{Special case}: $y_t = Ak_t$ (linear in capital, $\alpha = 1$, no labor).
	\vspace{-6pt}
	\begin{itemize}
		\item Unbounded ``endogenous'' growth
		\item Two countries starting with different $k_0$ remain forever different (percentage gap constant)
		\item Empirically implausible: labor earns $\sim$2/3 of income
	\end{itemize}
	
\end{frame}

\begin{frame}{Other Dynamics: Poverty Traps and Chaos}
	Poverty traps in the Solow model
		\vspace{-5pt}
\begin{figure}[h]
	\centering
	\includegraphics[scale=0.42]{FigsChapter3/Solow2R}
\end{figure}
\vspace{-13pt}	
	\begin{itemize}
		\item $\bar{k}_1$ and $\bar{k}_3$: stable; $\bar{k}_2$: unstable
		\item Economy starting with low $k_0$ converges to $\bar{k}_1$ (poverty trap)
		\item To escape: temporarily raise $s$ enough to push past $\bar{k}_2$
	\end{itemize}
	
\end{frame}
	
	\begin{frame}{Other Dynamics: Complex Dynamics}
		Complex dynamics in the Solow model
				\vspace{-5pt}
		\begin{figure}[h]
			\centering
			\includegraphics[scale=0.42]{FigsChapter3/Solow4R}
		\end{figure}
		\vspace{-13pt}	
	\begin{itemize}
		\item Locally stable if slope at steady state is $< 1$ in absolute value
		\item These cases are more esoteric in macroeconomics but illustrate the range of possible dynamics
	\end{itemize}
\end{frame}

% ============================================================
\section{The Growing Economy}
% ============================================================

\begin{frame}{Adding Technology and Population Growth}
		\vspace{-10pt}
	\[
	Y_t = F(K_t, A_t L_t)
	\]
	\vspace{-20pt}
	\begin{itemize}
		\item $A_t$: labor-augmenting (Harrod-neutral) technology; $A_t L_t$: effective labor
		\item Growth rates: $A_{t+1}/A_t = 1+\gamma$, $L_{t+1}/L_t = 1+n$
		\item Uzawa (1961): labor-augmenting is the \emph{only} form of technical change consistent with exact balanced growth 
	\end{itemize}
	

	Define capital per effective worker: $\tilde{k}_t \equiv K_t/(A_t L_t)$.
	

	Fundamental equation becomes:
	\[
	(1+\gamma)(1+n)\tilde{k}_{t+1} = (1-\delta)\tilde{k}_t + sf(\tilde{k}_t)
	\]
	
	Very similar to the basic model. Population growth $n$ and technology growth $\gamma$ act like additional ``depreciation'' on $\tilde{k}$.
\end{frame}

\begin{frame}{Balanced Growth Path}
	\textbf{Balanced growth path (BGP)}: $\tilde{k}_t$ is constant at $\bar{\tilde{k}}$:
	\[
	(1+\gamma)(1+n)\bar{\tilde{k}} = (1-\delta)\bar{\tilde{k}} + sf(\bar{\tilde{k}})
	\]
	
	\textbf{Cobb-Douglas}:
	\[
	\bar{\tilde{k}} = \left(\frac{s}{(1+\gamma)(1+n) + \delta - 1}\right)^{\frac{1}{1-\alpha}}
	\]
	
	\bigskip
	\textbf{Convergence}: Starting from any $\tilde{k}_0$, the sequence $\{\tilde{k}_t\}$ converges monotonically to $\bar{\tilde{k}}$.
	
	\bigskip
	\textbf{On BGP}: $Y_t/(A_t L_t) = f(\bar{\tilde{k}})$ is constant, so income per capita $y_t = Y_t/L_t = f(\bar{\tilde{k}})A_t$ grows at rate $\gamma$.
\end{frame}


\begin{frame}{Long-Run Growth}
	\textbf{Per capita income} on the BGP:
	\[
	y_t = f(\bar{\tilde{k}})A_t \qquad \Rightarrow \qquad \frac{y_{t+1} - y_t}{y_t} = \gamma
	\]
		\vspace{-20pt}
	\begin{itemize}
		\item \textbf{Long-run per capita growth rate $= \gamma$ (technology growth)}
		\item \textbf{No other parameters} affect the long-run growth rate
		\item Higher $s$ raises the \emph{level} of income, not the growth rate
		\item Higher $s$ also raises the growth rate in the \emph{short run} (during transition)
	\end{itemize}
	

	\textbf{Short-run growth}:
	\[
	\frac{y_{t+1} - y_t}{y_t} = \frac{f(\tilde{k}_{t+1})}{f(\tilde{k}_t)}(1+\gamma) - 1
	\]
		\vspace{-15pt}
	\begin{itemize}
		\item If $\tilde{k}_t < \bar{\tilde{k}}$: short-run growth $> \gamma$ (poor countries grow faster)
		\item If $\tilde{k}_t > \bar{\tilde{k}}$: short-run growth $< \gamma$ (rich countries grow slower)
	\end{itemize}
\end{frame}

% ============================================================
\section{Stylized Facts and the Solow Model}
% ============================================================

\begin{frame}{Matching the Kaldor Facts}
	\textbf{Constant $K/Y$}: On BGP, $K_t/(A_t L_t) = \bar{\tilde{k}}$ and $Y_t/(A_t L_t) = f(\bar{\tilde{k}})$ are both constant, so $K_t/Y_t$ is constant. $\checkmark$
	
	\bigskip
	\textbf{Steady per capita growth}: $y_t$ grows at rate $\gamma$ on BGP. $\checkmark$
	
	\bigskip
	\textbf{Constant return to capital}: Firms maximize profit under competition:
	\[
	\max_{K_t, A_t L_t} F(K_t, A_t L_t) - r_t K_t - w_t A_t L_t
	\]
	FOC: $r_t = F_1(K_t, A_t L_t) = f'(\tilde{k}_t)$. On BGP: $r_t = f'(\bar{\tilde{k}})$, which is constant. $\checkmark$
\end{frame}

\begin{frame}{Factor Shares}
	\textbf{Wage per efficiency unit}: From the FOC for $A_t L_t$:
	\[
	w_t = f(\tilde{k}_t) - \tilde{k}_t f'(\tilde{k}_t)
	\]
	Constant on BGP (since $\tilde{k}_t = \bar{\tilde{k}}$).
	
	\bigskip
	\textbf{Capital share}: $r_t K_t / Y_t = \tilde{k}_t f'(\tilde{k}_t)/f(\tilde{k}_t)$. Constant on BGP. $\checkmark$
	
	\bigskip
	\textbf{Labor share}: $w_t A_t L_t / Y_t = 1 - $ capital share (from CRS $\Rightarrow$ zero pure profits). Constant on BGP. $\checkmark$
	
	\bigskip
	\textbf{Cobb-Douglas}: $Y_t = K_t^\alpha (A_t L_t)^{1-\alpha}$
	\begin{itemize}
		\item Capital share $= \alpha$, labor share $= 1-\alpha$, regardless of $K_t$ or $A_t L_t$
		\item Factor shares constant even \emph{off} the BGP
	\end{itemize}
\end{frame}


% ============================================================
\section{Convergence}
% ============================================================

\begin{frame}{Speed of Convergence: Basic Model}
	Linearize the fundamental equation around $\bar{k}$:
	\[
	\Delta k_{t+1} = (1 - \delta + sf'(\bar{k}))\Delta k_t
	\]
	where $\Delta k_t \equiv k_t - \bar{k}$.
	

	Equivalently:
	\[
	\frac{k_{t+1} - \bar{k}}{\bar{k}} = (1-\lambda)\frac{k_t - \bar{k}}{\bar{k}}
	\]
	
	where $\lambda \equiv \delta - sf'(\bar{k})$ is the \textbf{convergence speed}.
	

	\textbf{Cobb-Douglas}: $\lambda = \delta(1-\alpha)$.
	\vspace{-6pt}
	\begin{itemize}
		\item Faster convergence when $\alpha$ is small (stronger decreasing returns)
		\item Faster convergence when $\delta$ is large
		\item $s$ does \emph{not} affect $\lambda$ in the Cobb-Douglas case (two opposing forces exactly cancel)
	\end{itemize}
\end{frame}

\begin{frame}{Speed of Convergence: Growing Economy}
	Slow and fast convergence
			\vspace{-5pt}
		\begin{columns}[T] % T = align at top
			
			\begin{column}{0.48\textwidth}
				\centering
				\includegraphics[width=\textwidth]{FigsChapter3/Solow5R}
	
			\end{column}
			
			\begin{column}{0.48\textwidth}
				\centering
				\includegraphics[width=\textwidth]{FigsChapter3/Solow6R}

			\end{column}
			
		\end{columns}
	Higher $\lambda$ = flatter slope at steady state = fewer steps to converge.\\
Linearizing the growing-economy fundamental equation:
	\[
	\frac{\tilde{k}_{t+1} - \bar{\tilde{k}}}{\bar{\tilde{k}}} = (1-\lambda)\frac{\tilde{k}_t - \bar{\tilde{k}}}{\bar{\tilde{k}}}
	\]
	\textbf{Cobb-Douglas}:
	$\displaystyle\lambda = (1-\alpha)\left(1 - \frac{1-\delta}{(1+\gamma)(1+n)}\right)$
	
\end{frame}



\begin{frame}{Cross-Country Evidence: All Countries}
	 All countries, 1960--2019
	\vspace{-5pt}
\begin{figure}[h]
	\centering
	\includegraphics[scale=0.42]{FigsChapter3/convergence1}
\end{figure}
\vspace{-13pt}	
	\begin{itemize}
		\item \textbf{No systematic tendency} for poor countries to grow faster
		\item Does not reject the Solow model: model predicts \textbf{conditional convergence} (countries converge if they share the same parameters), not \emph{unconditional} convergence
		\item Saving rates, $\gamma$, and institutions differ widely across countries
	\end{itemize}
\end{frame}

\begin{frame}{Cross-Country Evidence: OECD Countries}
	OECD countries, 1960--2019
		\vspace{-5pt}
	\begin{figure}[h]
		\centering
		\includegraphics[scale=0.42]{FigsChapter3/convergence2}
	\end{figure}
	\vspace{-13pt}	
	\begin{itemize}
		\item OECD: high-income countries with similar institutions 
		\item \textbf{Clear tendency for convergence}: initially poorer OECD countries grew faster
		\item Barro and Sala-i-Martin (1995): similar patterns within US states, Japanese prefectures, European regions
	\end{itemize}
\end{frame}

\begin{frame}{Recent Evidence: Unconditional Convergence Emerging}
	 All countries, 2000--2019
	\vspace{-5pt}
\begin{figure}[h]
	\centering
	\includegraphics[scale=0.42]{FigsChapter3/convergence3}
\end{figure}
\vspace{-13pt}	
	\begin{itemize}
		\item Kremer, Willis, and You (2022): in recent years, data show a tendency for \textbf{unconditional} convergence
		\item Negative correlation between GDP and subsequent growth
		\item Explanation: underlying factors (policies, institutions, human capital) have become more similar across countries
	\end{itemize}
\end{frame}

% ============================================================
\section{Calibration}
% ============================================================

\begin{frame}{Calibrating the Solow Model}
	\textbf{Goal}: Assign functional forms and parameter values to obtain quantitative predictions.
	
	\bigskip
	\textbf{Production function}: Cobb-Douglas $F(K,AL) = K^\alpha(AL)^{1-\alpha}$
	\begin{itemize}
		\item Consistent with approximately constant factor shares
		\item Substitution elasticity $\approx 1$ in aggregate production function studies
	\end{itemize}
	
	\bigskip
	\textbf{Parameters} (annual frequency):
	\begin{itemize}
		\item $\alpha = 1/3$ (capital share from NIPA, Figure 2.12)
		\item $\gamma = 0.02$ (long-run per capita income growth)
		\item $n = 0.01$ (population growth)
		\item $\delta = 0.046$ (implied by $I/K \approx 0.076$ and the BGP condition $(1+\gamma)(1+n) = 1 - \delta + I/K$)
	\end{itemize}
\end{frame}

\begin{frame}{Quantitative Predictions: Convergence Speed}
	With $\alpha = 1/3$, $\gamma = 0.02$, $n = 0.01$, $\delta = 0.046$:
	\[
	\lambda = (1-\alpha)\left(1 - \frac{1-\delta}{(1+\gamma)(1+n)}\right) = 0.049
	\]
	
	\begin{itemize}
		\item \textbf{Empirical estimate}: $\lambda \approx 0.015$ to $0.03$ (Barro and Sala-i-Martin, 2004)
		\item \textbf{Model over-predicts} convergence speed
		\item To match $\lambda \approx 0.02$: need $\alpha \approx 0.73$
		\item Justification for larger $\alpha$: part of labor income is return to \textbf{human capital} (accumulable, like physical capital); $A_t$ itself may be accumulable (R\&D investment)
	\end{itemize}
\end{frame}


\begin{frame}{Quantitative Experiments: Doubling the Saving Rate}
	With $s = 0.1$ and calibrated parameters:
	\begin{itemize}
		\item $\bar{\tilde{k}} = 1.20$, $f(\bar{\tilde{k}}) = 1.06$
	\end{itemize}
	
	\bigskip
	Doubling to $s = 0.2$:
	\begin{itemize}
		\item $\bar{\tilde{k}} = 1.90$, $f(\bar{\tilde{k}}) = 1.24$
	\end{itemize}
	
	\bigskip
	\textbf{Result}: Doubling the saving rate increases the normalized level of output by \textbf{17\%} ($1.24/1.06 = 1.17$).
	
	\bigskip
	Can also compute quantitative predictions for how fast output rises to reach this new level (transition dynamics).
\end{frame}

% ============================================================
\section{Business Cycles}
% ============================================================

\begin{frame}{The Solow Model as the Core of Business Cycle Theory}
	The Solow model's ability to account for long-run facts makes it the \textbf{core of macroeconomic modeling}. Virtually all modern business cycle theories build on it.
	
	\bigskip
	Business cycles $=$ arrival of shocks $+$ propagation mechanism.
	
	\bigskip
	\textbf{Three types of shocks} (extending the basic model):
	\begin{enumerate}
		\item Neutral technology shocks (RBC model)
		\item Investment-specific technology shocks
		\item Demand shocks (Keynesian)
	\end{enumerate}
	
	\bigskip
	In each case, $k_{t+1}$ can be written as a function of $k_t$ and the shock---a modified fundamental equation.
\end{frame}

\begin{frame}{RBC Model: Technology Shocks}
	Production function with variable labor and TFP shock:
	\[
	y_t = A_t F(k_t, \ell_t)
	\]
	With endogenous saving and labor supply (to match business cycle facts):
	\[
	k_{t+1} = (1-\delta)k_t + s(k_t, A_t) A_t F(k_t, \ell(k_t, A_t))
	\]
	Business cycle facts require:
	\vspace{-6pt}
	\begin{itemize}
		\item $\ell_t$ comoves positively with cycle
		\item Investment $i_t$ more volatile than $y_t$
		\item Consumption $c_t$ less volatile than $y_t$
	\end{itemize}
	
	These cannot be matched with constant $s$ and fixed $\ell$---need endogenous responses to $A_t$.
	
	\textbf{Note}: With Cobb-Douglas, Hicks-neutral ($A_t F(k,\ell)$) and Harrod-neutral ($F(k, A_t \ell)$) are equivalent: $K^\alpha(A\ell)^{1-\alpha} = A^{1-\alpha}K^\alpha\ell^{1-\alpha}$.
\end{frame}

\begin{frame}{Investment-Specific and Demand Shocks}
	\textbf{Investment-specific shocks}: Goods market becomes $y_t = c_t + i_t/\nu_t$, where $\nu_t$ varies. When $\nu_t$ is high, investment goods are cheaper:
	\[
	k_{t+1} = s\nu_t F(k_t, \ell) + (1-\delta)k_t
	\]
	
	\bigskip
	\textbf{Demand shocks} (Keynesian): Consumption $c_t$ is exogenous, $i_t/c_t = s/(1-s)$. When demand is insufficient, unemployment $u_t$ arises:
	\[
	y_t = \frac{1}{1-s}c_t = F(k_t, \ell(1-u_t))
	\]
	\[
	k_{t+1} = sF(k_t, \ell(1-u(c_t, k_t))) + (1-\delta)k_t
	\]
	
	This is Keynesian in spirit but requires frictions to explain why output falls below full capacity. Several chapters address these frictions.
\end{frame}

\begin{frame}{Impulse Response Functions}
	Impulse response---how $A_t$ (left) affects $k_t$ (right)
				\vspace{-5pt}
	\begin{columns}[T] % T = align at top
		
		\begin{column}{0.48\textwidth}
			\centering
			\includegraphics[width=\textwidth]{FigsChapter3/IRAR}
			
		\end{column}
		
		\begin{column}{0.48\textwidth}
			\centering
			\includegraphics[width=\textwidth]{FigsChapter3/IR1R}
			
		\end{column}
		
	\end{columns}
\vspace{-6pt}
	Experiment: $A_0 = (1+\varepsilon)\bar{A}$, $A_t = (1+\rho^t\varepsilon)\bar{A}$ for $t \geq 1$ in the baseline model
	
	\textbf{Three key features}:
	\vspace{-6pt}
	\begin{enumerate}
		\item \textbf{Slow adjustment}: $k_t$ response is far more persistent than $A_t$
		\item \textbf{Small magnitude}: maximum deviation of $k_t$ is only 0.25\%
		\item \textbf{$k_t$ returns to steady state}: convergence force remains active
	\end{enumerate}
\end{frame}

\begin{frame}{Log-Linearization}
	Express variables as percent deviations from steady state: $\hat{x}_t \equiv \log(x_t/\bar{x})$.
	
	\medskip
	Log-linearizing the fundamental equation around $\bar{k}$:
	\[
	\hat{k}_{t+1} = (1-\delta(1-\alpha))\hat{k}_t + \delta\hat{A}_t = (1-\lambda)\hat{k}_t + \delta\hat{A}_t
	\]
	
	With $\hat{A}_t = \rho^t\varepsilon$, the solution is:
	\[
	\hat{k}_{t+1} = \rho^t \cdot \frac{1 - \left(\frac{1-\lambda}{\rho}\right)^{t+1}}{1 - \frac{1-\lambda}{\rho}} \cdot \delta\varepsilon
	\]
	
	Log-linearized output response: $\hat{y}_t = \hat{A}_t + \alpha\hat{k}_t$.
	
	\medskip
	\textbf{Accuracy}: Maximum error $\approx 0.00001$ in levels---visually indistinguishable from the nonlinear solution.
\end{frame}

% ============================================================
\section{Summary}
% ============================================================

\begin{frame}{Key Takeaways (I)}
	\begin{enumerate}
		\item \textbf{Fundamental equation}: $k_{t+1} = (1-\delta)k_t + sf(k_t)$; unique steady state $\bar{k}$; global monotonic convergence
		
		\item \textbf{Convergence force}: Decreasing returns to capital---when $k$ is low, growth is fast; when $k$ is high, growth is slow
		
		\item \textbf{Growing economy}: Labor-augmenting technology is necessary for balanced growth (Uzawa); $\tilde{k}_t \to \bar{\tilde{k}}$; per capita growth rate $= \gamma$
		
		\item \textbf{Stylized facts}: Matches constant $K/Y$, constant factor shares, constant $r$, rising wages, steady per capita growth
	\end{enumerate}
\end{frame}

\begin{frame}{Key Takeaways (II)}
	\begin{enumerate}
		\setcounter{enumi}{4}
		\item \textbf{Convergence speed}: $\lambda = (1-\alpha)(1 - (1-\delta)/((1+\gamma)(1+n)))$; model over-predicts unless $\alpha$ is raised (human capital)
		
		\item \textbf{Cross-country data}: No unconditional convergence, but clear conditional convergence (OECD, US states); recent evidence of unconditional convergence post-2000
		
		\item \textbf{Business cycles}: Solow model is the backbone; technology, investment-specific, and demand shocks; impulse responses; log-linearization
		
		\item \textbf{Limitation}: Saving rate $s$ and labor $\ell$ are exogenous; Chapter 4 will endogenize them via optimization
	\end{enumerate}
\end{frame}

\end{document}

